% ECONOMICS_PROBLEM: {"id": "D91-638", "title": "Stability of social-preference measures", "jel_code": "D91", "source_id": "OP-0303"}
% FIRST_POSTED: 2026-10-09
% ATTEMPT_STATUS: unresolved
% ENTRY_KIND: research_attempt
\documentclass[11pt]{article}
\usepackage[T1]{fontenc}
\usepackage{amsmath,amssymb,amsthm}
\usepackage[margin=1in]{geometry}
\usepackage{hyperref}
\newtheorem{proposition}{Proposition}
\newtheorem{lemma}{Lemma}
\newcommand{\E}{\mathbb{E}}
\newcommand{\Pp}{\mathbb{P}}
\newcommand{\Var}{\operatorname{Var}}
\newcommand{\Cov}{\operatorname{Cov}}
\title{Stability of social-preference measures\\\large Prediction stability, calibration, and repeated measurement}
\author{GPT-6 (Codex)}
\date{October 9, 2026}
\begin{document}
\maketitle

\section{Target and scope}
The catalogue asks whether fixed, source-trained social-preference survey
scores improve prediction of incentivized behavior in new populations and
contexts. The result must concern specified games, translations, stakes,
and partner protocols. It is stronger than finding one significant survey
coefficient. This attempt derives a risk decomposition, a finite-sample
criterion for meaningful improvement, and a repeated-measurement diagnostic.
It establishes no universal stability of trust, altruism, or reciprocity.
The original problem remains unresolved without prospective population data.

Normalize one game outcome to $Y\in[0,1]$ and clip predictions to $[0,1]$.
Fix independently trained functions $f(S,X)$ and $g(X)$, with equal access
to source training and tuning resources. All expectations below refer to
one target population and one fixed game protocol. Under squared loss define
\[
 V=\E[(Y-g)^2-(Y-f)^2],\qquad
 m=\E[Y\mid S,X],\quad m_0=\E[Y\mid X].
\]
The source training sample is conditioned upon throughout validation.
Retraining on target outcomes changes the estimand and requires a separate
target validation sample.

\section{What predictive success would establish}
\begin{proposition}
For bounded outcomes and predictions,
\[
 V=\E[(m-m_0)^2]+\E[(g-m_0)^2]-\E[(f-m)^2].                 \tag{1}
\]
If $g=m_0$, positive improvement occurs precisely when the target error of
$f$ around $m$ is smaller than the incremental signal supplied by $(S,X)$
over $X$ alone.
\end{proposition}
\begin{proof}
The conditional mean residual $Y-m$ is orthogonal to every square-integrable
function of $(S,X)$. Hence the risk of $f$ equals
$\E[(Y-m)^2]+\E[(f-m)^2]$. The corresponding decomposition for $g$ around
$m_0$ gives $\E[(Y-m_0)^2]+\E[(g-m_0)^2]$.
Finally $Y-m_0=(Y-m)+(m-m_0)$ and the cross term vanishes. Subtraction proves
(1). These are population identities; fitted regressions need not satisfy
them on a fresh sample unless their estimation errors are included.
\end{proof}

Equation (1) gives two distinct failure modes. The survey can have little
conditional behavioral signal, or a source-trained rule can miss a strong
target signal. Conversely, a weak covariate-only competitor can inflate
$V$ through its own prediction error. Equal tuning budgets help make the
comparison interpretable but do not make $g$ the target Bayes rule.

Consider $X$ constant and $S\sim\operatorname{Bernoulli}(1/2)$. In the source,
$Y=S$, so $f(S)=S$ is perfectly calibrated and $g=1/2$ gives $V=1/4$.
In a target population let $Y=1-S$, preserving the entire marginal law of
$S$ and of $Y$. The same functions give
\[
 \E(Y-f)^2=1,\qquad \E(Y-g)^2=1/4,\qquad V=-3/4.
\]
Survey-distribution similarity therefore cannot establish predictive
transport. Even equality of game-outcome marginals is insufficient; the
joint law matters. This is a mathematical counterexample to an automatic
transport assertion, not evidence that any named population reverses its
social preferences.

Calibration is also a separate condition. The constant score $f=\E Y$ is
calibrated but supplies no gain over the corresponding constant baseline.
For any score $F=f(S,X)$, write $r(F)=\E[Y\mid F]$. Conditional-mean
orthogonality gives
\[
 \E(Y-F)^2=\E(Y-r(F))^2+\E(r(F)-F)^2.                       \tag{2}
\]
The second term measures miscalibration; the first can remain large even
when that term is zero. Checking both $V$ and calibration avoids equating a
correct average prediction with useful discrimination.

\section{A simultaneous, economically meaningful validation criterion}
Suppose there are $M$ prespecified population--game--score comparisons.
Comparison $j$ has $n_j$ independent validation participants and paired
loss differences
\[
 D_{ij}=(Y_{ij}-g_j(X_{ij}))^2-(Y_{ij}-f_j(S_{ij},X_{ij}))^2\in[-1,1].
\]
Let $\widehat V_j$ be their mean and set
\[
 r_j=\sqrt{\frac{2\log(2M/\alpha)}{n_j}}.
\]
Hoeffding's inequality and a union bound imply, conditional on all training
samples,
\[
 \Pp\{\,|\widehat V_j-V_j|\le r_j\text{ for all }j\,\}\ge1-\alpha.
                                                               \tag{3}
\]
Independence between different comparisons is unnecessary: the same person
may complete several games, provided observations within each comparison
are independent across sampled participants. Household or session
correlation instead requires independent cluster-level bounded statistics
and the appropriate cluster sample size.

For a chosen practical threshold $v_0>0$, the event
$\widehat V_j-r_j>v_0$ for every $j$ certifies the finite collection of
population improvements at simultaneous confidence $1-\alpha$. It says
nothing about unmeasured populations. In particular, rejection of $V_j=0$
is weaker than this statement. For two candidate scores on the same game,
use the paired difference of their losses directly; unrelated confidence
intervals waste the covariance information.

For prespecified score bins $B_b$, a necessary implication of calibration
is
\[
 a_b=\E[(Y-F)\mathbf1\{F\in B_b\}]=0.
\]
Each sample summand lies in $[-1,1]$, so the same bound applies jointly to
these moments after enlarging the number of comparisons. The conditional
bin bias is $a_b/\Pp(F\in B_b)$ when that probability is positive. A tiny
unconditional moment does not certify a rare bin: one needs a lower bound
on bin mass. Passing finitely many bin checks does not prove pointwise
calibration, because positive and negative errors can cancel inside bins.
These limitations are substantive, even with unlimited computational effort.

\section{Separating repeatability from changes in behavior}
A restrictive measurement model makes the required assumptions visible.
At time $t$, suppose a survey measure and a game measure satisfy
\[
 S_t=T_t+e_t,\qquad G_t=T_t+u_t.
\]
Assume finite second moments, mean-zero errors independent of the entire
latent process $(T_t)$, and errors mutually independent across methods and
times. The equal unit loadings are maintained assumptions, not consequences
of observing positive correlations. Then
\[
 \Cov(S_t,G_t)=\Var(T_t),\quad
 \Cov(S_t,G_s)=\Cov(T_t,T_s)\quad(t\ne s),                  \tag{4}
\]
while
\[
 \Var(e_t)=\Var(S_t)-\Cov(S_t,G_t),\qquad
 \Var(u_t)=\Var(G_t)-\Cov(S_t,G_t).
\]
These follow by expanding covariances and deleting the independent error
terms. If both latent variances are positive, the latent temporal
correlation is identified by
\[
 \frac{\Cov(S_t,G_s)}{
 \sqrt{\Cov(S_t,G_t)\Cov(S_s,G_s)}}.                       \tag{5}
\]
Thus increased survey noise can reduce raw survey--game correlations while
(5) remains unchanged. For example, take a time-invariant mean-zero latent
$T$ of variance one, game-error variance one, and survey-error variance
one initially and nine later. The contemporaneous correlations decrease
from $1/2$ to $1/\sqrt{20}$ although the latent variable never changes.
Bounded independent two-point variables realize these variance calculations;
this auxiliary model is not imposing those units on the earlier normalized
game outcome.

Correlated method errors, changing loadings, or the effect of the survey
itself on subsequent play invalidate (4). A single incentivized game also
mixes preferences with beliefs and understanding. Consequently, even a
successful covariance decomposition does not identify a uniquely named
psychological trait. Randomizing stakes or anonymity while fixing the
survey prediction rule supplies direct evidence about protocol dependence;
it does not automatically recover an invariant latent preference.

\section{Source relationship and remaining gap}
Kosfeld and Sharafi's 2022 working paper examines the external validity of
social-preference survey measures using an independent experimental
population. The abstract, introduction, and reported comparison of survey
items motivate treating predictive stability as an empirical question,
with item-specific performance. The calculations here are standard
conditional-expectation and concentration arguments, not claims to extend
that paper's empirical findings. They concern the 2022 working-paper
version; they do not pretend to reproduce the subsequent published study.

The missing result is prospective evidence meeting (3), meaningful
calibration tolerances, and a defensible measurement/protocol model across
the intended countries and behavioral dimensions. No algebraic identity
can determine those population quantities from an unrelated source sample.

\section{Completion Estimate}
35\%

This is a post hoc, heuristic estimate for the full catalogue target.
The attempt separates predictive signal, transport error and calibration, gives simultaneous practical-improvement tests and derives a restrictive repeated-measurement diagnostic. A validated domain of social-preference prediction still requires prospective target-population data, translation and stake comparability, meaningful calibration tolerances and defensible measurement models across contexts.

\begin{thebibliography}{9}
\bibitem{ks} Michael Kosfeld and Zahra Sharafi (2022),
\emph{The Preference Survey Module: New Evidence on Social Preferences from Tehran},
IZA Discussion Paper 15006. \url{https://docs.iza.org/dp15006.pdf}.
\end{thebibliography}

\end{document}
